\chapter*{Glosario}
\addcontentsline{toc}{chapter}{Glosario}

\begin{description}
\item[Juego Inmutable] Principio arquitectónico donde el motor del juego no contiene lógica específica de agentes de inteligencia artificial, limitándose a exponer estado y recibir acciones.

\item[Arquitectura Plug-and-Play] Diseño modular que permite intercambiar agentes de IA sin modificar el núcleo del videojuego.

\item[Middleware] Capa intermedia encargada de recibir el estado del juego, procesarlo, enrutar la decisión del agente y registrar métricas experimentales.

\item[Función de Transición] Función determinista $T(s,a)=s'$ que define cómo evoluciona el estado del entorno tras aplicar una acción.

\item[Equivalencia Funcional] Propiedad mediante la cual dos implementaciones del entorno producen el mismo estado resultante ante las mismas entradas.

\item[Golden Logs] Trazas generadas desde el entorno Unity utilizadas como referencia para validar la equivalencia del simulador externo.

\item[Política Baseline] Política heurística determinista utilizada como referencia experimental.

\item[Política Paramétrica] Política cuyo comportamiento depende de un conjunto explícito de pesos ajustables.

\item[Aleatoriedad Acotada] Mecanismo que introduce variabilidad controlada mediante un margen relativo $\epsilon$ sin alterar el rendimiento competitivo principal.

\item[EvaluationMode] Modo experimental donde se fija la semilla del generador aleatorio para garantizar reproducibilidad.

\item[Winrate] Proporción de partidas ganadas respecto al total evaluado.

\item[Diferencial de Puntaje] Diferencia promedio de puntaje entre dos agentes evaluados.

\item[Latencia de Decisión] Tiempo requerido por un agente para seleccionar una acción en un turno.

\item[Reproducibilidad Experimental] Capacidad de replicar resultados bajo mismas condiciones iniciales y configuración de semillas.

\item[Imitation Learning] Enfoque de aprendizaje supervisado donde un modelo es entrenado para aproximar las decisiones de una política base.

\item[Behavior Cloning] Técnica de aprendizaje por imitación donde un modelo (estudiante) se entrena con pares estado-acción generados por otro modelo (profesor o \textit{teacher}).

\item[Destilación de Conocimiento] Proceso de transferir el comportamiento aprendido de un modelo complejo a uno más liviano, preservando la mayor fidelidad posible.

\item[Response Table] Tabla de mapeo que asocia cada cluster (perfil de oponente) con el especialista más efectivo para ese perfil.

\item[Multi-seed Real] Técnica de entrenamiento donde cada genoma o modelo se evalúa contra múltiples semillas simultáneamente, promediando fitness para prevenir sobreajuste a una secuencia determinista particular.

\item[Sistema Adaptativo] Combinación de un perfilador (KMeans), una tabla de respuesta y un conjunto de especialistas que permite seleccionar dinámicamente la mejor estrategia en tiempo de ejecución.

\item[Especialista] Modelo de IA entrenado específicamente contra un estilo de juego particular (por ejemplo, greedy, denial o spread).

\item[Generalista] Modelo de IA que presenta buen rendimiento contra múltiples estilos sin estar especializado en ninguno.

\item[Aprendizaje Online Pleno] Forma de adaptación donde se actualizan los pesos o parámetros del modelo durante la partida, utilizando las interacciones en curso como datos de entrenamiento. Requiere ejecutar algoritmos de optimización en tiempo real.

\item[Adaptación Online por Selección] Forma de adaptación donde no se actualizan los pesos del modelo en runtime. Se mantiene un catálogo de políticas pre-entrenadas y se selecciona la más adecuada en función del perfil del oponente, requiriendo únicamente inferencia.

\item[Brecha de Datos] Relación entre el volumen de datos requerido para entrenar un modelo funcional y el volumen disponible en una sesión real contra un jugador humano. Valores altos indican inviabilidad práctica del entrenamiento online.
\end{description}
\vfill